% configurações
\documentclass{article}

% português
%\usepackage[utf8]{inputenc}
%\usepackage[portuges,brazilian]{babel}
%\usepackage[T1]{fontenc}

% pacotes
%\usepackage{tikz}
%\usepackage[siunitx]{circuitikz}
\usepackage{pst-circ}
%\usepackage{amssymb,amsmath}
%\usepackage{hyperref}

% símbolos
\newcommand{\ohm}{\Omega}
\newcommand{\milli}{\mu}

\title{Clock}

\begin{document}

%\maketitle

\begin{pspicture}[showgrid=true](10,10)
  \pnode(1,10){A}
  \pnode(2,10){B}
  \pnode(6,10){C}
  \pnode(8,10){D}
  \pnode(2,8){E}
  \wire[arrows=o-*](A)(B)
  \wire[arrows=*-*](B)(C)
  \wire[arrows=*-*](C)(D)
  \resistor[dipolestyle=zigzag](B)(E){$R1$}

\begin{pspicture}
  \logicic[nicpins=8,bubblesize=0.1,%
    pintl=true,pintllabel=tl,pintlnumber=1,%
    pintc=true,pintclabel=tc,pintcnumber=2,%
    pintr=true,pintrlabel=tr,pintrnumber=3,%
    invertpintl=true,invertpintc=true,invertpintr=true,%
    pinbl=true,pinbllabel=bl,pinblnumber=1,%
    pinbc=true,pinbclabel=bc,pinbcnumber=2,%
    pinbr=true,pinbrlabel=br,pinbrnumber=3,%
    invertpinbl=true,invertpinbc=true,invertpinbr=true,%
    pinalabel=a,pinblabel=b,pinclabel=c,pindlabel=d,%
    pinelabel=e,pinflabel=f,pinglabel=g,pinhlabel=h,%
    pinanumber=1,pinbnumber=2,pincnumber=3,pindnumber=4,%
    pinenumber=5,pinfnumber=6,pingnumber=7,pinhnumber=8](3,7){Name}
\end{pspicture}

\end{pspicture}

\end{document}
